% Options for packages loaded elsewhere
\PassOptionsToPackage{unicode}{hyperref}
\PassOptionsToPackage{hyphens}{url}
\PassOptionsToPackage{space}{xeCJK}
\documentclass[
  a4paper,
]{article}
\usepackage{xcolor}
\usepackage[margin=24mm]{geometry}
\usepackage{amsmath,amssymb}
\setcounter{secnumdepth}{-\maxdimen} % remove section numbering
\usepackage{iftex}
\ifPDFTeX
  \usepackage[T1]{fontenc}
  \usepackage[utf8]{inputenc}
  \usepackage{textcomp} % provide euro and other symbols
\else % if luatex or xetex
  \usepackage{unicode-math} % this also loads fontspec
  \defaultfontfeatures{Scale=MatchLowercase}
  \defaultfontfeatures[\rmfamily]{Ligatures=TeX,Scale=1}
\fi
\usepackage{lmodern}
\ifPDFTeX\else
  % xetex/luatex font selection
  \setmonofont[]{DejaVuSansMono.ttf}
  \ifXeTeX
    \usepackage{xeCJK}
    \setCJKmainfont[]{HaranoAjiMincho-Regular.otf}
  \fi
  \ifLuaTeX
    \usepackage[]{luatexja-fontspec}
    \setmainjfont[]{HaranoAjiMincho-Regular.otf}
  \fi
\fi
% Use upquote if available, for straight quotes in verbatim environments
\IfFileExists{upquote.sty}{\usepackage{upquote}}{}
\IfFileExists{microtype.sty}{% use microtype if available
  \usepackage[]{microtype}
  \UseMicrotypeSet[protrusion]{basicmath} % disable protrusion for tt fonts
}{}
\makeatletter
\@ifundefined{KOMAClassName}{% if non-KOMA class
  \IfFileExists{parskip.sty}{%
    \usepackage{parskip}
  }{% else
    \setlength{\parindent}{0pt}
    \setlength{\parskip}{6pt plus 2pt minus 1pt}}
}{% if KOMA class
  \KOMAoptions{parskip=half}}
\makeatother
\usepackage{color}
\usepackage{fancyvrb}
\newcommand{\VerbBar}{|}
\newcommand{\VERB}{\Verb[commandchars=\\\{\}]}
\DefineVerbatimEnvironment{Highlighting}{Verbatim}{commandchars=\\\{\}}
% Add ',fontsize=\small' for more characters per line
\newenvironment{Shaded}{}{}
\newcommand{\AlertTok}[1]{\textcolor[rgb]{1.00,0.00,0.00}{\textbf{#1}}}
\newcommand{\AnnotationTok}[1]{\textcolor[rgb]{0.38,0.63,0.69}{\textbf{\textit{#1}}}}
\newcommand{\AttributeTok}[1]{\textcolor[rgb]{0.49,0.56,0.16}{#1}}
\newcommand{\BaseNTok}[1]{\textcolor[rgb]{0.25,0.63,0.44}{#1}}
\newcommand{\BuiltInTok}[1]{\textcolor[rgb]{0.00,0.50,0.00}{#1}}
\newcommand{\CharTok}[1]{\textcolor[rgb]{0.25,0.44,0.63}{#1}}
\newcommand{\CommentTok}[1]{\textcolor[rgb]{0.38,0.63,0.69}{\textit{#1}}}
\newcommand{\CommentVarTok}[1]{\textcolor[rgb]{0.38,0.63,0.69}{\textbf{\textit{#1}}}}
\newcommand{\ConstantTok}[1]{\textcolor[rgb]{0.53,0.00,0.00}{#1}}
\newcommand{\ControlFlowTok}[1]{\textcolor[rgb]{0.00,0.44,0.13}{\textbf{#1}}}
\newcommand{\DataTypeTok}[1]{\textcolor[rgb]{0.56,0.13,0.00}{#1}}
\newcommand{\DecValTok}[1]{\textcolor[rgb]{0.25,0.63,0.44}{#1}}
\newcommand{\DocumentationTok}[1]{\textcolor[rgb]{0.73,0.13,0.13}{\textit{#1}}}
\newcommand{\ErrorTok}[1]{\textcolor[rgb]{1.00,0.00,0.00}{\textbf{#1}}}
\newcommand{\ExtensionTok}[1]{#1}
\newcommand{\FloatTok}[1]{\textcolor[rgb]{0.25,0.63,0.44}{#1}}
\newcommand{\FunctionTok}[1]{\textcolor[rgb]{0.02,0.16,0.49}{#1}}
\newcommand{\ImportTok}[1]{\textcolor[rgb]{0.00,0.50,0.00}{\textbf{#1}}}
\newcommand{\InformationTok}[1]{\textcolor[rgb]{0.38,0.63,0.69}{\textbf{\textit{#1}}}}
\newcommand{\KeywordTok}[1]{\textcolor[rgb]{0.00,0.44,0.13}{\textbf{#1}}}
\newcommand{\NormalTok}[1]{#1}
\newcommand{\OperatorTok}[1]{\textcolor[rgb]{0.40,0.40,0.40}{#1}}
\newcommand{\OtherTok}[1]{\textcolor[rgb]{0.00,0.44,0.13}{#1}}
\newcommand{\PreprocessorTok}[1]{\textcolor[rgb]{0.74,0.48,0.00}{#1}}
\newcommand{\RegionMarkerTok}[1]{#1}
\newcommand{\SpecialCharTok}[1]{\textcolor[rgb]{0.25,0.44,0.63}{#1}}
\newcommand{\SpecialStringTok}[1]{\textcolor[rgb]{0.73,0.40,0.53}{#1}}
\newcommand{\StringTok}[1]{\textcolor[rgb]{0.25,0.44,0.63}{#1}}
\newcommand{\VariableTok}[1]{\textcolor[rgb]{0.10,0.09,0.49}{#1}}
\newcommand{\VerbatimStringTok}[1]{\textcolor[rgb]{0.25,0.44,0.63}{#1}}
\newcommand{\WarningTok}[1]{\textcolor[rgb]{0.38,0.63,0.69}{\textbf{\textit{#1}}}}
\usepackage{longtable,booktabs,array}
\newcounter{none} % for unnumbered tables
\usepackage{calc} % for calculating minipage widths
% Correct order of tables after \paragraph or \subparagraph
\usepackage{etoolbox}
\makeatletter
\patchcmd\longtable{\par}{\if@noskipsec\mbox{}\fi\par}{}{}
\makeatother
% Allow footnotes in longtable head/foot
\IfFileExists{footnotehyper.sty}{\usepackage{footnotehyper}}{\usepackage{footnote}}
\makesavenoteenv{longtable}
\setlength{\emergencystretch}{3em} % prevent overfull lines
\providecommand{\tightlist}{%
  \setlength{\itemsep}{0pt}\setlength{\parskip}{0pt}}
\usepackage{bookmark}
\IfFileExists{xurl.sty}{\usepackage{xurl}}{} % add URL line breaks if available
\urlstyle{same}
\hypersetup{
  hidelinks,
  pdfcreator={LaTeX via pandoc}}

\author{}
\date{}

\begin{document}

\section{Curvature Renormalization and Perfect-Reflection Stability by
Curved Closed Stationary Waves
v2}\label{curvature-renormalization-and-perfect-reflection-stability-by-curved-closed-stationary-waves-v2}

\textbf{Subtitle:} Numerical constructive experiment on the detection,
re-selection, and exchange-interference reflection recovery of
curvature-induced relative phase leakage under the all-positive zero
closure \texttt{Σx\_n\^{}2=0} \textbf{Date:} 2026-07-11 \textbf{Author:}
Noriaki Kihara \textbf{Position:} Additional paper in the Wave
Information Readout series \textbf{Version DOI:} 10.5281/zenodo.21332874
\textbf{Concept DOI:} 10.5281/zenodo.21304039

V2 recalculates the same conditions after changing the formula for the
fermion-like reflection map.

\begin{center}\rule{0.5\linewidth}{0.5pt}\end{center}

\subsection{Abstract}\label{abstract}

This paper numerically constructs how an odd-harmonic complex wave
satisfying the all-positive zero closure

\begin{Shaded}
\begin{Highlighting}[]
\NormalTok{Q(x)=\textbackslash{}sum\_n x\_n\^{}2=0}
\end{Highlighting}
\end{Shaded}

is stabilized inside a curved local cell.

This closure condition is not introduced here as an auxiliary assumption
for the curvature test. It is Axiom 1 of the existing basic axiom system
of this paper series, together with namelessness and non-trivial
existence. In that system, it was already placed as the existence
condition for non-trivial complex phase waves. The present paper applies
Axiom 1 to curved local cells and exchange-interference reflection.

In the preceding paper, a local exchange-interference map was
constructed from a fermionic inverse-phase core, and a
direction-reversal output corresponding to complete elastic reflection
was obtained without using an external \texttt{q\ -\textgreater{}\ -q}
instruction. The present paper asks whether that exchange-interference
map breaks under curvature-induced relative phase leakage, and whether
it is recovered by re-selection into a closed stationary wave.

The curvature effect is not treated as a quantitative prediction of real
spacetime curvature. Instead, curvature-induced relative phase leakage
inside a local cell is modeled as

\begin{Shaded}
\begin{Highlighting}[]
\NormalTok{\textbackslash{}delta\_\{K,m\}.}
\end{Highlighting}
\end{Shaded}

If this leakage is inserted into a closure pair

\begin{Shaded}
\begin{Highlighting}[]
\NormalTok{x\_m,\textbackslash{}qquad ix\_m,}
\end{Highlighting}
\end{Shaded}

then

\begin{Shaded}
\begin{Highlighting}[]
\NormalTok{x\_m\^{}2+\textbackslash{}left(e\^{}\{i\textbackslash{}delta\_\{K,m\}\}ix\_m\textbackslash{}right)\^{}2}
\NormalTok{=}
\NormalTok{x\_m\^{}2\textbackslash{}left(1{-}e\^{}\{i2\textbackslash{}delta\_\{K,m\}\}\textbackslash{}right),}
\end{Highlighting}
\end{Shaded}

and a closure residual appears in general. If an internal phase
re-selection \texttt{\textbackslash{}beta\_\{K,m\}} satisfies

\begin{Shaded}
\begin{Highlighting}[]
\NormalTok{\textbackslash{}delta\_\{K,m\}+\textbackslash{}beta\_\{K,m\}=0,}
\end{Highlighting}
\end{Shaded}

the closure condition is recovered.

In the minimal experiment, the transient state with curvature-induced
relative phase leakage produced closure-pair RMS
\texttt{1.2319416790092972e-02} and transmission leakage
\texttt{1.1503183254481797e-01}. After internal phase re-selection, the
stationary state recovered closure-pair RMS
\texttt{9.4283259783636047e-19} and transmission leakage \texttt{0.0}.

In the broad verification, eight curvature-relative-phase models were
tested against seven correction freedoms: \texttt{none},
\texttt{constant}, \texttt{linear}, \texttt{affine}, \texttt{quadratic},
\texttt{cubic}, and \texttt{full}. At maximum curvature relative phase
\texttt{1.2}, the uncorrected case left maximum transmission leakage
\texttt{1.6202719613622976e-01}, whereas the \texttt{full} correction
recovered closure-pair RMS \texttt{7.8949412793793227e-19} and
transmission leakage \texttt{0.0}.

In the integrated one-sided scattering verification, the residual
curvature phase obtained from the broad sweep was returned to the local
exchange-interference map. The uncorrected case produced maximum dynamic
transmission leakage \texttt{1.6202719613622971e-01}, whereas the
\texttt{full} correction reduced it to \texttt{1.6608667985580024e-19}.
The dynamic scattering matched the two-channel expectation with maximum
error \texttt{5.551115123125783e-17}, and the maximum norm error was
\texttt{4.440892098500626e-16}.

Within the numerical constructive scope of this paper, curvature effects
are therefore not absent. They appear in transient states as closure
residuals and transmission leakage. However, once the system is
re-selected into a closed stationary wave satisfying
\texttt{Σx\_n\^{}2=0}, the curvature-induced relative phase leakage is
absorbed into the internal phase configuration, and the
complete-reflection readout is recovered.

\textbf{Keywords:} all-positive zero closure, curved closed stationary
wave, curvature renormalization, odd harmonics, exchange interference,
complete reflection, local constant-curvature cell, internal phase
re-selection

\begin{center}\rule{0.5\linewidth}{0.5pt}\end{center}

\subsection{1. Introduction}\label{introduction}

\subsubsection{1.1 Background}\label{background}

This series adopts namelessness, all-positive zero closure, and
non-trivial existence as its basic axioms. In particular, all-positive
zero closure is Axiom 1 of the basic axiom system v2. It is not an
auxiliary assumption added for the present paper. The central closure
condition is

\begin{Shaded}
\begin{Highlighting}[]
\NormalTok{\textbackslash{}sum\_n x\_n\^{}2=0.}
\end{Highlighting}
\end{Shaded}

This is not the conjugate norm

\begin{Shaded}
\begin{Highlighting}[]
\NormalTok{\textbackslash{}sum\_n |x\_n|\^{}2.}
\end{Highlighting}
\end{Shaded}

It is a closure condition in which each component is squared as it is
and summed with all-positive signs.

In the minimal two-component case,

\begin{Shaded}
\begin{Highlighting}[]
\NormalTok{A\^{}2+(iA)\^{}2=0,}
\end{Highlighting}
\end{Shaded}

so sign reversal is generated internally without using an external
negative coefficient. Complex phase is therefore not an ornament; it
appears as the phase algebra required for non-trivial closure.

The preceding paper, ``Interference Construction of a Perfect Reflection
Map from a Fermionic Inverse-Phase Core,'' constructed a local
exchange-interference map from an internal inverse-phase core and two
exchange paths. It generated a direction-reversal readout corresponding
to complete reflection without using an external
\texttt{q\ -\textgreater{}\ -q} instruction.

The present paper investigates whether this exchange-interference
reflection remains stable when placed in a curved local cell and exposed
to curvature-induced relative phase leakage.

\subsubsection{1.2 Question}\label{question}

The question of this paper is:

\begin{quote}
When curvature-induced relative phase leakage occurs in a curved local
cell, can the all-positive zero closure \texttt{Σx\_n\^{}2=0} detect the
leakage and recover the complete-reflection readout through internal
phase re-selection into a closed stationary wave?
\end{quote}

The important point is that curvature effects are not ignored at the
outset. Instead, this paper first confirms that curvature-induced
relative phase leakage produces closure residuals and transmission
leakage, and then tests whether those residuals disappear after
re-selection into a closed stationary wave.

\begin{center}\rule{0.5\linewidth}{0.5pt}\end{center}

\subsection{2. What This Paper Does Not
Claim}\label{what-this-paper-does-not-claim}

This paper does not claim the following.

{\def\LTcaptype{none} % do not increment counter
\begin{longtable}[]{@{}
  >{\raggedright\arraybackslash}p{(\linewidth - 2\tabcolsep) * \real{0.5000}}
  >{\raggedright\arraybackslash}p{(\linewidth - 2\tabcolsep) * \real{0.5000}}@{}}
\toprule\noalign{}
\begin{minipage}[b]{\linewidth}\raggedright
Not claimed
\end{minipage} & \begin{minipage}[b]{\linewidth}\raggedright
Reason
\end{minipage} \\
\midrule\noalign{}
\endhead
\bottomrule\noalign{}
\endlastfoot
Quantitative prediction of real spacetime curvature effects & No
standard spacetime metric, field equation, or experimental unit system
is used \\
Curvature is always unobservable & Closure residuals and transmission
leakage appear in transient states \\
A unique physical formula for the curvature relative phase \texttt{δ\_K}
& It is used here as a local relative-phase model for verification \\
Derivation of standard quantum theory or general relativity & This is a
numerical constructive experiment on an internal axiom system \\
Derivation of a continuous interaction Hamiltonian & This paper treats
local maps and existence conditions of closed stationary waves \\
\end{longtable}
}

The claim of this paper is a minimal numerical construction: a closed
stationary wave satisfying \texttt{Σx\_n\^{}2=0} detects
curvature-induced relative phase leakage and recovers the
exchange-interference reflection readout through internal phase
re-selection.

\begin{center}\rule{0.5\linewidth}{0.5pt}\end{center}

\subsection{3. Basic Axiom and the Closure Null
Cone}\label{basic-axiom-and-the-closure-null-cone}

\subsubsection{3.1 All-Positive Zero
Closure}\label{all-positive-zero-closure}

This section restates Axiom 1 of the basic axiom system v2. The present
paper uses this axiom as its starting point, but it is not introduced
here ad hoc to make curvature renormalization work.

Axiom 1 sets the closure condition as

\begin{Shaded}
\begin{Highlighting}[]
\NormalTok{Q(x)=\textbackslash{}sum\_\{n=1\}\^{}N x\_n\^{}2=0.}
\end{Highlighting}
\end{Shaded}

The set of closed states is denoted by

\begin{Shaded}
\begin{Highlighting}[]
\NormalTok{\textbackslash{}mathcal N}
\NormalTok{=}
\NormalTok{\textbackslash{}\{x\textbackslash{}mid Q(x)=0\textbackslash{}\}.}
\end{Highlighting}
\end{Shaded}

\texttt{\textbackslash{}mathcal\ N} is not the set of states with zero
conjugate norm. It is the complex closure null cone defined by the
all-positive sum of squares.

\subsubsection{3.2 Closure Pair}\label{closure-pair}

The minimal closure pair is

\begin{Shaded}
\begin{Highlighting}[]
\NormalTok{x\_m,}
\NormalTok{\textbackslash{}qquad}
\NormalTok{ix\_m.}
\end{Highlighting}
\end{Shaded}

Then

\begin{Shaded}
\begin{Highlighting}[]
\NormalTok{x\_m\^{}2+(ix\_m)\^{}2=0.}
\end{Highlighting}
\end{Shaded}

For an odd-harmonic complex wave, \texttt{m} is the harmonic label and

\begin{Shaded}
\begin{Highlighting}[]
\NormalTok{h\_m=2m+1}
\end{Highlighting}
\end{Shaded}

is used.

\begin{center}\rule{0.5\linewidth}{0.5pt}\end{center}

\subsection{4. Curved Local Cell and Relative Phase
Leakage}\label{curved-local-cell-and-relative-phase-leakage}

\subsubsection{4.1 Two Types of Curvature
Action}\label{two-types-of-curvature-action}

Curvature action is divided into two types.

The first is common-factor action:

\begin{Shaded}
\begin{Highlighting}[]
\NormalTok{x\_m\textbackslash{}mapsto g\_\{K,m\}x\_m,}
\NormalTok{\textbackslash{}qquad}
\NormalTok{ix\_m\textbackslash{}mapsto g\_\{K,m\}ix\_m.}
\end{Highlighting}
\end{Shaded}

In this case,

\begin{Shaded}
\begin{Highlighting}[]
\NormalTok{(g\_\{K,m\}x\_m)\^{}2+(g\_\{K,m\}ix\_m)\^{}2=0,}
\end{Highlighting}
\end{Shaded}

and the closure pair is preserved.

The second is relative-phase action:

\begin{Shaded}
\begin{Highlighting}[]
\NormalTok{x\_m\textbackslash{}mapsto x\_m,}
\NormalTok{\textbackslash{}qquad}
\NormalTok{ix\_m\textbackslash{}mapsto e\^{}\{i\textbackslash{}delta\_\{K,m\}\}ix\_m.}
\end{Highlighting}
\end{Shaded}

In this case,

\begin{Shaded}
\begin{Highlighting}[]
\NormalTok{x\_m\^{}2+(e\^{}\{i\textbackslash{}delta\_\{K,m\}\}ix\_m)\^{}2}
\NormalTok{=}
\NormalTok{x\_m\^{}2\textbackslash{}left(1{-}e\^{}\{i2\textbackslash{}delta\_\{K,m\}\}\textbackslash{}right),}
\end{Highlighting}
\end{Shaded}

and a closure residual appears in general.

\subsubsection{4.2 Closed Stationary Wave}\label{closed-stationary-wave}

A wave that remains stable in a curved local cell is called a closed
stationary wave \texttt{x\_K}.

It satisfies

\begin{Shaded}
\begin{Highlighting}[]
\NormalTok{Q(x\_K)=0}
\end{Highlighting}
\end{Shaded}

and

\begin{Shaded}
\begin{Highlighting}[]
\NormalTok{\textbackslash{}mathcal U\_K x\_K=e\^{}\{i\textbackslash{}alpha\_K\}x\_K.}
\end{Highlighting}
\end{Shaded}

Here \texttt{\textbackslash{}mathcal\ U\_K} is the local evolution map
after passing through the curved local cell, and
\texttt{e\^{}\{i\textbackslash{}alpha\_K\}} is a common phase removed by
external readout.

\begin{center}\rule{0.5\linewidth}{0.5pt}\end{center}

\subsection{5. Internal Re-Selection of Curvature Relative
Phase}\label{internal-re-selection-of-curvature-relative-phase}

For relative phase leakage \texttt{\textbackslash{}delta\_\{K,m\}},
introduce an internal phase re-selection
\texttt{\textbackslash{}beta\_\{K,m\}}.

If

\begin{Shaded}
\begin{Highlighting}[]
\NormalTok{\textbackslash{}delta\_\{K,m\}+\textbackslash{}beta\_\{K,m\}=0}
\end{Highlighting}
\end{Shaded}

holds, then

\begin{Shaded}
\begin{Highlighting}[]
\NormalTok{e\^{}\{i(\textbackslash{}delta\_\{K,m\}+\textbackslash{}beta\_\{K,m\})\}=1,}
\end{Highlighting}
\end{Shaded}

and the closure pair is recovered.

Thus curvature-induced relative phase leakage is not erased. It is
re-selected into an internal phase configuration that can exist as a
closed stationary wave.

\begin{center}\rule{0.5\linewidth}{0.5pt}\end{center}

\subsection{6. Connection to Exchange-Interference
Reflection}\label{connection-to-exchange-interference-reflection}

In the preceding exchange-interference reflection paper, the effective
phase was

\begin{Shaded}
\begin{Highlighting}[]
\NormalTok{\textbackslash{}Delta\_\{\textbackslash{}mathrm\{eff\}\}=\textbackslash{}Delta\_F,}
\end{Highlighting}
\end{Shaded}

and the pure inverse-phase core
\texttt{\textbackslash{}Delta\_F=\textbackslash{}pi} produced complete
reflection.

In this paper, including curvature-induced relative phase leakage,

\begin{Shaded}
\begin{Highlighting}[]
\NormalTok{\textbackslash{}Delta\_\{\textbackslash{}mathrm\{eff\},m\}}
\NormalTok{=}
\NormalTok{\textbackslash{}Delta\_F+\textbackslash{}delta\_\{K,m\}+\textbackslash{}beta\_\{K,m\}.}
\end{Highlighting}
\end{Shaded}

The complete-reflection condition is

\begin{Shaded}
\begin{Highlighting}[]
\NormalTok{\textbackslash{}Delta\_\{\textbackslash{}mathrm\{eff\},m\}=\textbackslash{}pi.}
\end{Highlighting}
\end{Shaded}

Without curvature correction,

\begin{Shaded}
\begin{Highlighting}[]
\NormalTok{\textbackslash{}Delta\_\{\textbackslash{}mathrm\{eff\},m\}=\textbackslash{}pi+\textbackslash{}delta\_\{K,m\},}
\end{Highlighting}
\end{Shaded}

so transmission leakage appears.

If

\begin{Shaded}
\begin{Highlighting}[]
\NormalTok{\textbackslash{}beta\_\{K,m\}={-}\textbackslash{}delta\_\{K,m\},}
\end{Highlighting}
\end{Shaded}

then

\begin{Shaded}
\begin{Highlighting}[]
\NormalTok{\textbackslash{}Delta\_\{\textbackslash{}mathrm\{eff\},m\}=\textbackslash{}pi}
\end{Highlighting}
\end{Shaded}

is recovered, and the complete-reflection readout returns.

\begin{center}\rule{0.5\linewidth}{0.5pt}\end{center}

\subsection{7. Numerical Experiments}\label{numerical-experiments}

\subsubsection{7.1 Experiment 1: Minimal Closed Stationary Wave
Verification}\label{experiment-1-minimal-closed-stationary-wave-verification}

The executed script is:

\begin{Shaded}
\begin{Highlighting}[]
\NormalTok{run\_curved\_closure\_stationary\_wave\_v2.py}
\end{Highlighting}
\end{Shaded}

The output directory is:

\begin{Shaded}
\begin{Highlighting}[]
\NormalTok{curved\_closure\_stationary\_wave\_result\_v2/}
\end{Highlighting}
\end{Shaded}

The curvature relative phase leakage is set as

\begin{Shaded}
\begin{Highlighting}[]
\NormalTok{\textbackslash{}delta\_\{K,m\}=\textbackslash{}kappa h\_m.}
\end{Highlighting}
\end{Shaded}

For maximum \texttt{\textbackslash{}kappa=0.012} and maximum curvature
phase \texttt{1.1879999999999999}, the following values were obtained.

{\def\LTcaptype{none} % do not increment counter
\begin{longtable}[]{@{}lrr@{}}
\toprule\noalign{}
State & Closure-pair RMS & Transmission leakage \\
\midrule\noalign{}
\endhead
\bottomrule\noalign{}
\endlastfoot
flat & \texttt{0.0000000000000000e+00} &
\texttt{0.0000000000000000e+00} \\
conformal & \texttt{9.4283259783636047e-19} &
\texttt{0.0000000000000000e+00} \\
transient & \texttt{1.2319416790092972e-02} &
\texttt{1.1503183254481797e-01} \\
stationary & \texttt{9.4283259783636047e-19} &
\texttt{0.0000000000000000e+00} \\
\end{longtable}
}

In the relaxation sequence, the final closure-pair RMS was
\texttt{2.0285753500943141e-09}, and the final transmission leakage was
\texttt{2.4915322496409796e-15}.

All verdict flags were \texttt{true}.

\subsubsection{7.2 Experiment 2: Broad Sweep of Curvature Phase
Models}\label{experiment-2-broad-sweep-of-curvature-phase-models}

The executed script is:

\begin{Shaded}
\begin{Highlighting}[]
\NormalTok{run\_curved\_closure\_stationary\_wave\_broad\_sweep\_v2.py}
\end{Highlighting}
\end{Shaded}

The output directory is:

\begin{Shaded}
\begin{Highlighting}[]
\NormalTok{curved\_closure\_stationary\_wave\_broad\_sweep\_result\_v2/}
\end{Highlighting}
\end{Shaded}

The eight curvature relative phase models were:

\begin{Shaded}
\begin{Highlighting}[]
\NormalTok{linear}
\NormalTok{quadratic\_area}
\NormalTok{cubic\_high}
\NormalTok{quartic\_edge}
\NormalTok{alternating\_linear}
\NormalTok{sinusoidal\_loop}
\NormalTok{mixed\_smooth}
\NormalTok{rippled\_random\_like}
\end{Highlighting}
\end{Shaded}

The seven internal correction freedoms were:

\begin{Shaded}
\begin{Highlighting}[]
\NormalTok{none}
\NormalTok{constant}
\NormalTok{linear}
\NormalTok{affine}
\NormalTok{quadratic}
\NormalTok{cubic}
\NormalTok{full}
\end{Highlighting}
\end{Shaded}

The correction-wise aggregate at maximum curvature relative phase
\texttt{1.2} was:

{\def\LTcaptype{none} % do not increment counter
\begin{longtable}[]{@{}
  >{\raggedright\arraybackslash}p{(\linewidth - 6\tabcolsep) * \real{0.2000}}
  >{\raggedleft\arraybackslash}p{(\linewidth - 6\tabcolsep) * \real{0.2667}}
  >{\raggedleft\arraybackslash}p{(\linewidth - 6\tabcolsep) * \real{0.2667}}
  >{\raggedleft\arraybackslash}p{(\linewidth - 6\tabcolsep) * \real{0.2667}}@{}}
\toprule\noalign{}
\begin{minipage}[b]{\linewidth}\raggedright
Correction
\end{minipage} & \begin{minipage}[b]{\linewidth}\raggedleft
Max closure-pair RMS
\end{minipage} & \begin{minipage}[b]{\linewidth}\raggedleft
Max transmission leakage
\end{minipage} & \begin{minipage}[b]{\linewidth}\raggedleft
Max residual phase
\end{minipage} \\
\midrule\noalign{}
\endhead
\bottomrule\noalign{}
\endlastfoot
none & \texttt{1.3992789770439718e-02} & \texttt{1.6202719613622976e-01}
& \texttt{1.2000000000000000e+00} \\
constant & \texttt{1.2323357129304359e-02} &
\texttt{1.1632365832184562e-01} & \texttt{1.1885985027360557e+00} \\
linear & \texttt{1.2317721991993871e-02} &
\texttt{1.1624067781252551e-01} & \texttt{1.2122605013795469e+00} \\
affine & \texttt{1.2315221532542749e-02} &
\texttt{1.1619972192010980e-01} & \texttt{1.2235371932078578e+00} \\
quadratic & \texttt{1.2300337923101595e-02} &
\texttt{1.1598359742226144e-01} & \texttt{1.2740926312051919e+00} \\
cubic & \texttt{1.2278927640523116e-02} &
\texttt{1.1567739316276086e-01} & \texttt{1.3311608455010591e+00} \\
full & \texttt{7.8949412793793227e-19} & \texttt{0.0000000000000000e+00}
& \texttt{0.0000000000000000e+00} \\
\end{longtable}
}

Common-factor curvature preserved closure throughout the sweep. Without
correction, non-trivial curvature leakage was detected. With
\texttt{full} correction, closure and complete reflection were recovered
for all phase models. With limited corrections, model-dependent
residuals remained.

\subsubsection{7.3 Experiment 3: Integrated Verification with One-Sided
Scattering}\label{experiment-3-integrated-verification-with-one-sided-scattering}

The executed script is:

\begin{Shaded}
\begin{Highlighting}[]
\NormalTok{run\_curved\_closure\_scattering\_integration\_v2.py}
\end{Highlighting}
\end{Shaded}

The output directory is:

\begin{Shaded}
\begin{Highlighting}[]
\NormalTok{curved\_closure\_scattering\_integration\_result\_v2/}
\end{Highlighting}
\end{Shaded}

The residual curvature relative phases from Experiment 2 were returned
to the local exchange-interference map for one-sided incident
scattering.

The correction-wise aggregate was:

{\def\LTcaptype{none} % do not increment counter
\begin{longtable}[]{@{}
  >{\raggedright\arraybackslash}p{(\linewidth - 6\tabcolsep) * \real{0.2000}}
  >{\raggedleft\arraybackslash}p{(\linewidth - 6\tabcolsep) * \real{0.2667}}
  >{\raggedleft\arraybackslash}p{(\linewidth - 6\tabcolsep) * \real{0.2667}}
  >{\raggedleft\arraybackslash}p{(\linewidth - 6\tabcolsep) * \real{0.2667}}@{}}
\toprule\noalign{}
\begin{minipage}[b]{\linewidth}\raggedright
Correction
\end{minipage} & \begin{minipage}[b]{\linewidth}\raggedleft
Max dynamic transmission leakage
\end{minipage} & \begin{minipage}[b]{\linewidth}\raggedleft
Max closure-pair RMS
\end{minipage} & \begin{minipage}[b]{\linewidth}\raggedleft
Max expectation error
\end{minipage} \\
\midrule\noalign{}
\endhead
\bottomrule\noalign{}
\endlastfoot
none & \texttt{1.6202719613622971e-01} & \texttt{1.3992789770439718e-02}
& \texttt{5.5511151231257827e-17} \\
constant & \texttt{1.1632365832184560e-01} &
\texttt{1.2323357129304359e-02} & \texttt{1.3877787807814457e-17} \\
linear & \texttt{1.1624067781252549e-01} &
\texttt{1.2317721991993871e-02} & \texttt{1.3877787807814457e-17} \\
affine & \texttt{1.1619972192010980e-01} &
\texttt{1.2315221532542749e-02} & \texttt{1.0408340855860843e-17} \\
quadratic & \texttt{1.1598359742226147e-01} &
\texttt{1.2300337923101595e-02} & \texttt{2.7755575615628914e-17} \\
cubic & \texttt{1.1567739316276089e-01} &
\texttt{1.2278927640523116e-02} & \texttt{2.7755575615628914e-17} \\
full & \texttt{1.6608667985580024e-19} & \texttt{7.8949412793793227e-19}
& \texttt{1.6608667985580024e-19} \\
\end{longtable}
}

The dynamic scattering matched the two-channel expectation with maximum
error \texttt{5.551115123125783e-17}, and the maximum norm error was
\texttt{4.440892098500626e-16}.

\begin{center}\rule{0.5\linewidth}{0.5pt}\end{center}

\subsection{8. Classification of
Results}\label{classification-of-results}

The results are classified as follows.

{\def\LTcaptype{none} % do not increment counter
\begin{longtable}[]{@{}
  >{\raggedright\arraybackslash}p{(\linewidth - 4\tabcolsep) * \real{0.3333}}
  >{\raggedright\arraybackslash}p{(\linewidth - 4\tabcolsep) * \real{0.3333}}
  >{\raggedright\arraybackslash}p{(\linewidth - 4\tabcolsep) * \real{0.3333}}@{}}
\toprule\noalign{}
\begin{minipage}[b]{\linewidth}\raggedright
Target
\end{minipage} & \begin{minipage}[b]{\linewidth}\raggedright
Classification
\end{minipage} & \begin{minipage}[b]{\linewidth}\raggedright
Verdict
\end{minipage} \\
\midrule\noalign{}
\endhead
\bottomrule\noalign{}
\endlastfoot
\texttt{Σx\_n\^{}2=0} requires non-trivial complex closure & Derived
consequence & retained \\
Curvature relative phase leakage produces closure residuals &
Numerically constructed consequence & retained \\
Common-factor curvature action does not break closure & Numerically
constructed consequence & retained \\
Internal phase re-selection recovers closure & Numerically constructed
consequence & retained \\
Internal phase re-selection recovers complete reflection & Numerically
constructed consequence & retained \\
Real-spacetime local flatness is explained by this structure &
Connection to existing theory & not claimed \\
\end{longtable}
}

\begin{center}\rule{0.5\linewidth}{0.5pt}\end{center}

\subsection{9. Discussion}\label{discussion}

\subsubsection{9.1 Curvature Has Not
Disappeared}\label{curvature-has-not-disappeared}

The result of this paper is not that curvature has no effect.

In the transient state with curvature-induced relative phase leakage,
closure residuals and transmission leakage clearly appear.

Curvature effects are therefore detectable in this model.

\subsubsection{9.2 Closed Stationary Waves Absorb Curvature
Internally}\label{closed-stationary-waves-absorb-curvature-internally}

The important case is the state after curvature-induced relative phase
leakage has occurred and the system is re-selected as a closed
stationary wave.

If the internal phase correction \texttt{\textbackslash{}beta\_\{K,m\}}
satisfies

\begin{Shaded}
\begin{Highlighting}[]
\NormalTok{\textbackslash{}delta\_\{K,m\}+\textbackslash{}beta\_\{K,m\}=0,}
\end{Highlighting}
\end{Shaded}

then closure residuals and transmission leakage disappear.

This does not mean that curvature is absent. It means that the curvature
relative phase is renormalized into the internal phase configuration of
the closed stationary wave that can stably exist.

\subsubsection{9.3 Relation to Local
Flatness}\label{relation-to-local-flatness}

In the standard local-flatness approximation, curvature effects are
ignored in a sufficiently small region.

The construction in this paper is different.

Curvature effects are first inserted, closure is observed to break, and
then the residual disappears from readout only after the state is
re-selected as a closed stationary wave.

Thus, within the internal axiom system of this paper, a readout that
looks locally flat can be interpreted not as the result of ignoring
curvature, but as the result that only curvature-inclusive stationary
interference modes that close are stably read out.

This is not an identification with standard theory. It is a working
hypothesis for constructing a correspondence map.

\subsubsection{9.4 Connection to Complete Elastic
Reflection}\label{connection-to-complete-elastic-reflection}

The complete-reflection map constructed in the preceding paper was an
exchange-interference map with an internal inverse-phase core:

\begin{Shaded}
\begin{Highlighting}[]
\NormalTok{\textbackslash{}Delta\_F=\textbackslash{}pi.}
\end{Highlighting}
\end{Shaded}

In the present paper, adding curvature-induced relative phase leakage
gives

\begin{Shaded}
\begin{Highlighting}[]
\NormalTok{\textbackslash{}Delta\_\{\textbackslash{}mathrm\{eff\},m\}}
\NormalTok{=}
\NormalTok{\textbackslash{}pi+\textbackslash{}delta\_\{K,m\}+\textbackslash{}beta\_\{K,m\}.}
\end{Highlighting}
\end{Shaded}

Without correction, \texttt{\textbackslash{}delta\_\{K,m\}} remains and
transmission leakage appears. Under \texttt{full} correction,
\texttt{\textbackslash{}beta\_\{K,m\}=-\textbackslash{}delta\_\{K,m\}},
so

\begin{Shaded}
\begin{Highlighting}[]
\NormalTok{\textbackslash{}Delta\_\{\textbackslash{}mathrm\{eff\},m\}=\textbackslash{}pi}
\end{Highlighting}
\end{Shaded}

is recovered.

At that point, the one-sided scattering experiment also recovers
\texttt{R=1,T=0}.

\begin{center}\rule{0.5\linewidth}{0.5pt}\end{center}

\subsection{10. Conclusion}\label{conclusion}

This paper numerically constructed how an odd-harmonic complex wave
satisfying the all-positive zero closure \texttt{Σx\_n\^{}2=0} is
stabilized inside a curved local cell.

In the minimal experiment, the transient state with curvature-induced
relative phase leakage produced closure-pair RMS
\texttt{1.2319416790092972e-02} and transmission leakage
\texttt{1.1503183254481797e-01}. After internal phase re-selection, the
stationary state recovered closure-pair RMS
\texttt{9.4283259783636047e-19} and transmission leakage \texttt{0.0}.

In the broad verification, eight curvature relative phase models and
seven correction freedoms were compared. Without correction, maximum
transmission leakage \texttt{1.6202719613622976e-01} remained. Under
\texttt{full} correction, closure-pair RMS
\texttt{7.8949412793793227e-19} and transmission leakage \texttt{0.0}
were recovered.

In the integrated one-sided local exchange-interference scattering
verification, the uncorrected case produced maximum dynamic transmission
leakage \texttt{1.6202719613622971e-01}, while the \texttt{full}
correction reduced maximum dynamic transmission leakage to
\texttt{1.6608667985580024e-19}.

Therefore, within the numerical constructive scope of this paper,
curvature effects are not ignored. Curvature-induced relative phase
leakage appears in transient states as closure residuals and
transmission leakage. However, once the system is re-selected into a
closed stationary wave satisfying \texttt{Σx\_n\^{}2=0}, the
curvature-induced relative phase leakage is absorbed into the internal
phase configuration, and the complete-reflection readout by exchange
interference is recovered.

Thus \texttt{Σx\_n\^{}2=0} functions not merely as a conservation
condition, but as an existence condition for non-trivial complex waves,
a detection condition for curvature-induced relative phase leakage, a
re-selection condition into closed stationary waves, and a stability
condition for the complete-reflection readout.

\begin{center}\rule{0.5\linewidth}{0.5pt}\end{center}

\section{Appendix A. Executed Programs and
Outputs}\label{appendix-a.-executed-programs-and-outputs}

\subsection{A.1 Minimal Verification}\label{a.1-minimal-verification}

\begin{Shaded}
\begin{Highlighting}[]
\NormalTok{python3 run\_curved\_closure\_stationary\_wave\_v2.py}
\end{Highlighting}
\end{Shaded}

Output:

\begin{Shaded}
\begin{Highlighting}[]
\NormalTok{curved\_closure\_stationary\_wave\_result\_v2/}
\end{Highlighting}
\end{Shaded}

Main files:

{\def\LTcaptype{none} % do not increment counter
\begin{longtable}[]{@{}
  >{\raggedright\arraybackslash}p{(\linewidth - 2\tabcolsep) * \real{0.5000}}
  >{\raggedright\arraybackslash}p{(\linewidth - 2\tabcolsep) * \real{0.5000}}@{}}
\toprule\noalign{}
\begin{minipage}[b]{\linewidth}\raggedright
Type
\end{minipage} & \begin{minipage}[b]{\linewidth}\raggedright
File
\end{minipage} \\
\midrule\noalign{}
\endhead
\bottomrule\noalign{}
\endlastfoot
Report &
\href{curved_closure_stationary_wave_result_v2/curved_closure_stationary_wave_report_v2.md}{curved\_closure\_stationary\_wave\_report\_v2.md} \\
JSON &
\href{curved_closure_stationary_wave_result_v2/curved_closure_stationary_wave_result_v2.json}{curved\_closure\_stationary\_wave\_result\_v2.json} \\
sweep CSV &
\href{curved_closure_stationary_wave_result_v2/curved_closure_stationary_wave_sweep_v2.csv}{curved\_closure\_stationary\_wave\_sweep\_v2.csv} \\
relaxation CSV &
\href{curved_closure_stationary_wave_result_v2/curved_closure_stationary_wave_relaxation_v2.csv}{curved\_closure\_stationary\_wave\_relaxation\_v2.csv} \\
sweep figure &
\href{curved_closure_stationary_wave_result_v2/curved_closure_stationary_wave_sweep_v2.png}{curved\_closure\_stationary\_wave\_sweep\_v2.png} \\
relaxation figure &
\href{curved_closure_stationary_wave_result_v2/curved_closure_stationary_wave_relaxation_v2.png}{curved\_closure\_stationary\_wave\_relaxation\_v2.png} \\
\end{longtable}
}

\subsection{A.2 Broad Verification}\label{a.2-broad-verification}

\begin{Shaded}
\begin{Highlighting}[]
\NormalTok{python3 run\_curved\_closure\_stationary\_wave\_broad\_sweep\_v2.py}
\end{Highlighting}
\end{Shaded}

Output:

\begin{Shaded}
\begin{Highlighting}[]
\NormalTok{curved\_closure\_stationary\_wave\_broad\_sweep\_result\_v2/}
\end{Highlighting}
\end{Shaded}

Main files:

{\def\LTcaptype{none} % do not increment counter
\begin{longtable}[]{@{}
  >{\raggedright\arraybackslash}p{(\linewidth - 2\tabcolsep) * \real{0.5000}}
  >{\raggedright\arraybackslash}p{(\linewidth - 2\tabcolsep) * \real{0.5000}}@{}}
\toprule\noalign{}
\begin{minipage}[b]{\linewidth}\raggedright
Type
\end{minipage} & \begin{minipage}[b]{\linewidth}\raggedright
File
\end{minipage} \\
\midrule\noalign{}
\endhead
\bottomrule\noalign{}
\endlastfoot
Report &
\href{curved_closure_stationary_wave_broad_sweep_result_v2/curved_closure_stationary_wave_broad_report_v2.md}{curved\_closure\_stationary\_wave\_broad\_report\_v2.md} \\
JSON &
\href{curved_closure_stationary_wave_broad_sweep_result_v2/curved_closure_stationary_wave_broad_sweep_result_v2.json}{curved\_closure\_stationary\_wave\_broad\_sweep\_result\_v2.json} \\
sweep CSV &
\href{curved_closure_stationary_wave_broad_sweep_result_v2/curved_closure_stationary_wave_broad_sweep_v2.csv}{curved\_closure\_stationary\_wave\_broad\_sweep\_v2.csv} \\
control CSV &
\href{curved_closure_stationary_wave_broad_sweep_result_v2/curved_closure_stationary_wave_broad_conformal_control_v2.csv}{curved\_closure\_stationary\_wave\_broad\_conformal\_control\_v2.csv} \\
aggregate figure &
\href{curved_closure_stationary_wave_broad_sweep_result_v2/curved_closure_stationary_wave_broad_aggregate_v2.png}{curved\_closure\_stationary\_wave\_broad\_aggregate\_v2.png} \\
closure heatmap &
\href{curved_closure_stationary_wave_broad_sweep_result_v2/curved_closure_stationary_wave_broad_closure_heatmap_v2.png}{curved\_closure\_stationary\_wave\_broad\_closure\_heatmap\_v2.png} \\
leakage heatmap &
\href{curved_closure_stationary_wave_broad_sweep_result_v2/curved_closure_stationary_wave_broad_leakage_heatmap_v2.png}{curved\_closure\_stationary\_wave\_broad\_leakage\_heatmap\_v2.png} \\
\end{longtable}
}

\subsection{A.3 One-Sided Scattering
Integration}\label{a.3-one-sided-scattering-integration}

\begin{Shaded}
\begin{Highlighting}[]
\NormalTok{python3 run\_curved\_closure\_scattering\_integration\_v2.py}
\end{Highlighting}
\end{Shaded}

Output:

\begin{Shaded}
\begin{Highlighting}[]
\NormalTok{curved\_closure\_scattering\_integration\_result\_v2/}
\end{Highlighting}
\end{Shaded}

Main files:

{\def\LTcaptype{none} % do not increment counter
\begin{longtable}[]{@{}
  >{\raggedright\arraybackslash}p{(\linewidth - 2\tabcolsep) * \real{0.5000}}
  >{\raggedright\arraybackslash}p{(\linewidth - 2\tabcolsep) * \real{0.5000}}@{}}
\toprule\noalign{}
\begin{minipage}[b]{\linewidth}\raggedright
Type
\end{minipage} & \begin{minipage}[b]{\linewidth}\raggedright
File
\end{minipage} \\
\midrule\noalign{}
\endhead
\bottomrule\noalign{}
\endlastfoot
Report &
\href{curved_closure_scattering_integration_result_v2/curved_closure_scattering_integration_report_v2.md}{curved\_closure\_scattering\_integration\_report\_v2.md} \\
JSON &
\href{curved_closure_scattering_integration_result_v2/curved_closure_scattering_integration_result_v2.json}{curved\_closure\_scattering\_integration\_result\_v2.json} \\
CSV &
\href{curved_closure_scattering_integration_result_v2/curved_closure_scattering_integration_v2.csv}{curved\_closure\_scattering\_integration\_v2.csv} \\
Figure &
\href{curved_closure_scattering_integration_result_v2/curved_closure_scattering_integration_v2.png}{curved\_closure\_scattering\_integration\_v2.png} \\
\end{longtable}
}

\begin{center}\rule{0.5\linewidth}{0.5pt}\end{center}

\section{References}\label{references}

\subsection{Self-Citations}\label{self-citations}

\begin{enumerate}
\def\labelenumi{\arabic{enumi}.}
\tightlist
\item
  Noriaki Kihara, ``Basic Axiom System v2 for the Nameless
  Equal-Amplitude Composite Wave Model,'' 2026-07-10.
\item
  Noriaki Kihara, ``Paper 0: Distortion of the Geodesic Unit Cell in
  Positively-Curved Constant-Curvature Space --- Exact Evaluation of
  Edge, Angle, Area, and Volume,'' Version DOI:
  \texttt{10.5281/zenodo.21303433}, Concept DOI:
  \texttt{10.5281/zenodo.20680269}, 2026.
\item
  Noriaki Kihara, ``Interference Construction of a Perfect Reflection
  Map from a Fermionic Inverse-Phase Core,'' Version DOI:
  \texttt{10.5281/zenodo.21332867}, Concept DOI:
  \texttt{10.5281/zenodo.21295479}, 2026.
\end{enumerate}

\subsection{External References}\label{external-references}

\begin{enumerate}
\def\labelenumi{\arabic{enumi}.}
\setcounter{enumi}{3}
\tightlist
\item
  H. S. M. Coxeter, \emph{Regular Polytopes}, 3rd ed., Dover, 1973.
\item
  S. Pancharatnam, ``Generalized theory of interference, and its
  applications,'' \emph{Proceedings of the Indian Academy of Sciences
  A}, 44, 247-262, 1956.
\item
  M. V. Berry, ``Quantal phase factors accompanying adiabatic changes,''
  \emph{Proceedings of the Royal Society of London A}, 392, 45-57, 1984.
  DOI: \texttt{10.1098/rspa.1984.0023}.
\end{enumerate}

\end{document}
